\section*{The Fabric Purification Sciences Complex at the U}

At the University of Utah, the new Fabric Purification Sciences Complex (the U's
ultra-serious name for a fancy campus laundromat) is running an experiment in
the basement of the MEB. Every student who swipes their UCard to start a wash
cycle generates a data point in a log.

For one day, the system records a sequence of integers. Each integer represents
the ``grime index'' of a load of laundry, in the exact order the machines
processed them.

Engineers at the lab have noticed that the system behaves most predictably
during stretches where this grime index never goes down from one load to the
next. In other words, starting from some load and moving forward:
\begin{quote}
As long as each new load is at least as dirty as the previous one, the lab
considers that part of the log a stable interval.
\end{quote}

The number of loads in such an interval is called its \emph{stability score}.

You can start such an interval at any position in the log.

When you hit a load that is less dirty than the previous one, that interval
ends.

Your job (before you head to Rice--Eccles for the game) is to help the
researchers figure out how ``stable'' the lab was that day.

They want to know the highest stability score over the entire log.